% ==========================================================================
%  Kompetenzen-Tool — Benutzerhandbuch
%  Kompilieren:  lualatex Benutzerhandbuch.tex  (zweimal für TOC)
% ==========================================================================
\documentclass[12pt,a4paper]{article}

\usepackage[top=2.5cm,bottom=2.5cm,left=2.8cm,right=2.8cm]{geometry}
\usepackage{fontspec}
\setmainfont{Latin Modern Sans}
\usepackage{amsmath}
\usepackage{xcolor}
\usepackage{tabularray}
\usepackage{graphicx}
\usepackage{microtype}
\usepackage{setspace}
\usepackage{hyperref}
\usepackage{enumitem}
\usepackage{titlesec}
\usepackage{fancyhdr}
\usepackage{mdframed}
\usepackage{luacode}

% ---------- Farben --------------------------------------------------------
\definecolor{primary}{RGB}{0,99,142}
\definecolor{accent}{RGB}{66,144,179}
\definecolor{lightbg}{RGB}{240,248,252}
\definecolor{warnbg}{RGB}{255,249,230}
\definecolor{warnborder}{RGB}{200,150,0}
\definecolor{tipbg}{RGB}{236,252,238}
\definecolor{tipborder}{RGB}{34,139,34}
\definecolor{adminbg}{RGB}{245,240,255}
\definecolor{adminborder}{RGB}{100,60,180}

% ---------- Seitenformatierung --------------------------------------------
\pagestyle{fancy}
\fancyhf{}
\fancyhead[L]{\small\color{primary}\textbf{Kompetenzen-Tool} — Benutzerhandbuch}
\fancyhead[R]{\small\color{gray}\thepage}
\fancyfoot[C]{\small\color{gray}Evangelische Gemeinschaftsschule Erfurt · \texttt{http://zeugnistool.tgsef.intern}}
\renewcommand{\headrulewidth}{0.4pt}
\renewcommand{\footrulewidth}{0.4pt}

% ---------- Überschriften -------------------------------------------------
\titleformat{\section}
  {\Large\bfseries\color{primary}}{\thesection}{1em}{}[\vspace{0.2em}\hrule\vspace{0.5em}]
\titleformat{\subsection}
  {\large\bfseries\color{primary!80!black}}{\thesubsection}{1em}{}
\titleformat{\subsubsection}
  {\normalsize\bfseries}{\thesubsubsection}{1em}{}

% ---------- Hilfsboxen ----------------------------------------------------
\newmdenv[
  backgroundcolor=lightbg,
  linecolor=primary,
  linewidth=1.5pt,
  leftline=true, topline=false, bottomline=false, rightline=false,
  innerleftmargin=10pt, innerrightmargin=10pt,
  innertopmargin=8pt, innerbottommargin=8pt,
]{infobox}

\newmdenv[
  backgroundcolor=warnbg,
  linecolor=warnborder,
  linewidth=1.5pt,
  leftline=true, topline=false, bottomline=false, rightline=false,
  innerleftmargin=10pt, innerrightmargin=10pt,
  innertopmargin=8pt, innerbottommargin=8pt,
]{warnbox}

\newmdenv[
  backgroundcolor=tipbg,
  linecolor=tipborder,
  linewidth=1.5pt,
  leftline=true, topline=false, bottomline=false, rightline=false,
  innerleftmargin=10pt, innerrightmargin=10pt,
  innertopmargin=8pt, innerbottommargin=8pt,
]{tipbox}

\newmdenv[
  backgroundcolor=adminbg,
  linecolor=adminborder,
  linewidth=2pt,
  leftline=true, topline=false, bottomline=false, rightline=false,
  innerleftmargin=10pt, innerrightmargin=10pt,
  innertopmargin=8pt, innerbottommargin=8pt,
]{adminbox}

\newcommand{\info}[1]{\begin{infobox}\textbf{Hinweis:} #1\end{infobox}}
\newcommand{\warn}[1]{\begin{warnbox}\textbf{Achtung:} #1\end{warnbox}}
\newcommand{\tip}[1]{\begin{tipbox}\textbf{Tipp:} #1\end{tipbox}}
\newcommand{\adminonly}{\begin{adminbox}\textbf{Admin-Bereich:}
Dieser Abschnitt richtet sich an Personen mit Admin-Zugang (Schulleitung, IT-Beauftrage).\end{adminbox}}

\newcommand{\menu}[1]{\textbf{\textcolor{primary}{#1}}}
\newcommand{\key}[1]{\texttt{\small[#1]}}
\newcommand{\field}[1]{\textit{„#1"}}

\setlength{\parindent}{0pt}
\setlength{\parskip}{6pt}

% ==========================================================================
\begin{document}

% ---------- Titelseite ----------------------------------------------------
\begin{titlepage}
  \centering
  \vspace*{3cm}
  {\Huge\bfseries\color{primary} Kompetenzen-Tool\par}
  \vspace{0.5cm}
  {\LARGE\color{primary!70!black} Benutzerhandbuch\par}
  \vspace{1.5cm}
  \hrule height 2pt
  \vspace{0.5cm}
  {\large Evangelische Gemeinschaftsschule Erfurt\par}
  \vspace{0.3cm}
  {\normalsize Digitales Zeugniserfassungs- und Exportwerkzeug\par}
  \vspace{0.5cm}
  \hrule height 0.5pt
  \vfill
  \begin{infobox}
    \textbf{Erreichbarkeit:}\\
    Nur im lokalen Schulnetz (WLAN oder VPN)\\[4pt]
    \texttt{http://zeugnistool.tgsef.intern}
  \end{infobox}
  \vspace{2cm}
  {\small\color{gray} Stand: \today}
\end{titlepage}

% ---------- Inhaltsverzeichnis --------------------------------------------
\tableofcontents
\newpage

% ==========================================================================
% TEIL I – FÜR LEHRKRÄFTE
% ==========================================================================

\section{Einleitung}

Das \textbf{Kompetenzen-Tool} unterstützt Lehrkräfte der Evangelischen Gemeinschaftsschule Erfurt
bei der Erstellung kompetenzorientierter Zeugnisse.
Im Browser werden Noten, Kompetenzen und Zeugnistexte erfasst — das fertige PDF wird automatisch erstellt.

\textbf{Das Tool übernimmt für Sie:}
\begin{itemize}[noitemsep]
  \item Kompetenzen pro Fach und Klasse auswählen
  \item Niveau und Themenurteile pro Schüler eingeben
  \item Persönliche Zeugnistexte mit Texteditor schreiben
  \item Bewertungen für Förderkinder (LB, GB) anpassen
  \item PDF-Zeugnisse automatisch erstellen
\end{itemize}

% ==========================================================================
\section{Zugang}

\begin{warnbox}
Das Tool ist \textbf{nur im lokalen Schulnetz} erreichbar.
Voraussetzung: Laptop im Schul-WLAN \textbf{oder} VPN-Verbindung zum Schulserver.
\end{warnbox}

\begin{center}
  \large\texttt{\textcolor{primary}{http://zeugnistool.tgsef.intern}}
\end{center}

Aktuelle Version von Firefox, Chrome, Edge oder Safari verwenden.

% ==========================================================================
\section{Anmelden}

Öffnen Sie die Adresse im Browser. Geben Sie \field{Benutzernamen} und \field{Passwort} ein.

\begin{tblr}{
  width=\linewidth,
  colspec={Q[l,wd=3cm] X[l]},
  hlines, vlines,
  cell{1}{1-2}={bg=primary,fg=white,font=\bfseries},
}
Rolle & Zugriff \\
Lehrer & Kompetenzen, Schülerdaten, Stammdaten \\
Admin & Alle Bereiche inkl. Setup, Übersicht, Export, Benutzerverwaltung \\
\end{tblr}

\tip{Passwörter können nur durch einen Admin geändert werden.
Wenden Sie sich bei Problemen an die IT-Beauftragten.}

% ==========================================================================
\section{Kompetenzen auswählen}

Für jede Klasse und jedes Fach legen Sie fest, welche Kompetenzen im Zeugnis erscheinen.
Nur ausgewählte Kompetenzen werden gedruckt.

\textbf{Schritt für Schritt:}
\begin{enumerate}[noitemsep]
  \item \menu{Kompetenzen} im Menü links öffnen
  \item Klasse, Fach und Block auswählen
  \item Themen aufklappen (auf den Pfeil klicken)
  \item Passende Kompetenzen anhaken
  \item \key{Speichern} klicken
\end{enumerate}

\info{Werkstätten: Alle Kompetenzen sind automatisch ausgewählt und können nicht geändert werden.}

\tip{Am Ende jedes Themas können eigene Kompetenzen ergänzt werden,
die nur für diese Klasse gelten.}

% ==========================================================================
\section{Förderkinder: LB und GB}

Bevor Sie Noten eintragen, sollten Sie wissen, wie LB- und GB-Schüler
im Tool behandelt werden.

\begin{tblr}{
  width=\linewidth,
  colspec={Q[l,wd=2.5cm] X[l] Q[l,wd=3.5cm]},
  hlines, vlines,
  cell{1}{1-3}={bg=primary,fg=white,font=\bfseries},
}
Markierung & Bewertungsform & Anzeige in Tabelle \\
Normal & Niveau 1–3, Noten 1–4 oder ne & Weiß \\
LB & Freitext statt Niveau und Noten & Grün \\
GB & Nur Freitext im Niveau-Feld, keine Noten & Orange \\
\end{tblr}

\subsection{LB-Schüler einrichten}
\begin{enumerate}[noitemsep]
  \item \menu{Stammdaten} öffnen, Klasse wählen
  \item Bei dem Schüler den Haken bei \field{LB} setzen
  \item Speichern
\end{enumerate}
Danach erscheinen in \menu{Schülerdaten} Freitext-Felder statt Noten-Dropdowns.

\subsection{Lebenspraxis}
Das Fach \textbf{Lebenspraxis} erscheint in \menu{Schülerdaten} automatisch,
sobald ein Schüler als LB oder GB markiert ist.
Es hat kein Kompetenz-Raster — nur ein Freitext-Feld für die Bewertung.

% ==========================================================================
\section{Schülerdaten: Noten und Niveau}

\menu{Schülerdaten} → Klasse und Fach auswählen.
Die Tabelle zeigt alle Schüler mit Niveau-Spalte und einer Spalte pro Thema.

\subsection{Niveau}

\begin{tblr}{
  width=\linewidth,
  colspec={Q[l,wd=2.5cm] X[l]},
  hlines, vlines,
  cell{1}{1-2}={bg=primary,fg=white,font=\bfseries},
}
Wert & Bedeutung \\
1 & Anforderungsebene 1 (grundlegend) \\
2 & Anforderungsebene 2 (Regelstandard) \\
3 & Anforderungsebene 3 (erweiternd) \\
9 & Sport: kein Niveau-Eintrag \\
ne & Nicht erteilt — ganzes Fach \\
— & Kein Niveau (Werkstätten, DuG, Mitarbeit) \\
\end{tblr}

\subsection{Themenurteile}
Note pro Thema: \textbf{1} (sehr gut) bis \textbf{4} (nicht ausreichend).\\
Sonderwert \textbf{ne} = nicht erteilt (das Thema erscheint im PDF als Blockeintrag).

\subsection{LB- und GB-Schüler in der Tabelle}
\begin{itemize}[noitemsep]
  \item \textbf{LB (grün):} Niveau-Feld und Noten sind Freitext-Eingaben.
    Beim Niveau öffnet sich ein Texteditor (Stift-Button).
  \item \textbf{GB (orange):} Nur das Niveau-Feld ist ein Freitext.
    Keine Themenurteile.
\end{itemize}

\tip{Nach dem Eintragen \key{Speichern} nicht vergessen.
Die Übersicht zeigt, welche Felder noch fehlen.}

% ==========================================================================
\section{Stammdaten und Zeugnistexte}

\menu{Stammdaten} → Klasse wählen.

\subsection{Was Sie hier bearbeiten}
\begin{itemize}[noitemsep]
  \item \textbf{Fehltage} und \textbf{Fehlstunden} (entschuldigt / unentschuldigt)
  \item \textbf{LB / GB} Markierung (beeinflusst Bewertungsform in Schülerdaten)
\end{itemize}

\subsection{Zeugnistext schreiben}
Klicken Sie auf den \key{Stift-Button} in der Zeugnistext-Spalte.
Ein Editor öffnet sich — geben Sie den persönlichen Beurteilungstext ein.

\subsubsection{Unterstützte Formatierungen}
\begin{tblr}{
  width=\linewidth,
  colspec={X[l] Q[l,wd=3cm]},
  hlines, vlines,
  cell{1}{1-2}={bg=primary,fg=white,font=\bfseries},
}
Formatierung & Tastenkürzel \\
Fett & Strg+B \\
Kursiv & Strg+I \\
Unterstrichen & Strg+U \\
Durchgestrichen & — \\
Aufzählung & — \\
Nummerierte Liste & — \\
Tabelle & — \\
Zentrierter Text & — \\
\end{tblr}

\warn{Sonderzeichen wie \$ \% \# \& \{ \} bitte vermeiden —
sie können das Zeugnis-PDF beschädigen.}

\tip{Der Text wird automatisch gespeichert, wenn Sie den Editor schließen.
Der Speichern-Button der Hauptseite gilt nur für Fehltage.}

\subsection{Bemerkungen}
Über den \key{Chat-Button} können Bemerkungen erfasst werden,
die ebenfalls im Zeugnis erscheinen können.

% ==========================================================================
\section{Übersicht: Bearbeitungsstand prüfen}

\menu{Übersicht} → Klasse wählen.
Hier sehen Sie auf einen Blick, was noch fehlt.

\subsection{Tab: Kompetenzen}
Zeigt für jedes Fach, wie viele Kompetenzen ausgewählt sind.
Fehlende Fächer sind rot markiert.

\subsection{Tab: Noten}

\begin{tblr}{
  width=\linewidth,
  colspec={Q[l,wd=3cm] X[l]},
  hlines, vlines,
  cell{1}{1-2}={bg=primary,fg=white,font=\bfseries},
}
Anzeige & Bedeutung \\
3/5 & 3 von 5 Themenurteilen eingetragen \\
\textcolor{green}{5/5} & Alle Noten vollständig \\
\textcolor{red}{0/5} & Noch keine Noten \\
✓ & Text vorhanden (LB/GB-Schüler) \\
✗ & Text fehlt noch \\
ZT ✓ & Zeugnistext vorhanden \\
\end{tblr}

\subsection{Tab: Eigene Kompetenzen}
Listet selbst hinzugefügte Kompetenzen dieser Klasse.
Können direkt hier bearbeitet werden.

% ==========================================================================
\section{Häufige Fragen}

\subsection*{Das Tool ist nicht erreichbar}
Prüfen Sie: WLAN verbunden? VPN aktiv?
Öffnen Sie \texttt{http://zeugnistool.tgsef.intern} im Browser.

\subsection*{Ich sehe keine Klassen}
Entweder wurde noch keine Datenbank für dieses Halbjahr ausgewählt,
oder noch keine Schüler importiert. Wenden Sie sich an einen Admin.

\subsection*{Ein LB-Schüler erscheint nicht in Lebenspraxis}
Das Fach erscheint nur, wenn der Schüler als LB oder GB markiert ist.
\menu{Stammdaten} → LB-Haken setzen → Speichern.

\subsection*{Der Zeugnistext verliert Formatierung}
Formatierungen wie Fett, Listen usw. bleiben erhalten, sofern sie
über den eingebauten Editor eingegeben wurden.
Bitte keinen Text aus Word einfügen — stattdessen zuerst in Editor eingeben.

\subsection*{Ich möchte eine alte Datenbank wiederherstellen}
\menu{Setup} → \menu{Backup} → gewünschte Sicherung wiederherstellen.
(Admin-Zugang erforderlich.)

\subsection*{Der Export schlägt fehl}
Häufigste Ursache: Zeugnistag nicht gesetzt.
Prüfen Sie \menu{Setup} → \menu{Zeugnistag}. (Admin-Zugang erforderlich.)

\newpage
% ==========================================================================
% TEIL II – FÜR ADMINISTRATOREN
% ==========================================================================

{%
  \centering
  \vspace{2em}
  {\color{adminborder}\rule{\linewidth}{2pt}}\par\vspace{0.5em}
  {\Large\bfseries\color{adminborder} Teil II: Administration}\par
  \vspace{0.3em}
  {\normalsize\color{gray} Für Schulleitung und IT-Beauftragte}\par
  \vspace{0.5em}
  {\color{adminborder}\rule{\linewidth}{2pt}}\par
  \vspace{2em}
}

% ==========================================================================
\section{Erste Schritte: Datenbank und Schuljahr einrichten}

\adminonly

\subsection{Neue Datenbank anlegen}
Pro Halbjahr oder Klassengruppe empfiehlt sich eine eigene Datenbank.

\begin{enumerate}[noitemsep]
  \item \menu{Setup} öffnen
  \item Im Abschnitt \field{Datenbanken} auf \key{Neue Datenbank} klicken
  \item Namen eingeben, z.\,B. \texttt{reports\_2026\_hj1}
  \item \key{Erstellen} bestätigen
  \item \menu{Schema} → \key{Schema initialisieren}
\end{enumerate}

\info{Empfohlene Namenskonvention: \texttt{reports\_[Jahr]\_[hj1|hj2|ej]}.
Keine Leerzeichen oder Sonderzeichen im Namen.}

\subsection{Zeugnistag setzen}
\menu{Setup} → \menu{Zeugnistag} → Datum im Format TT.MM.JJJJ eingeben.
Dieses Datum erscheint auf der Rückseite aller Zeugnisse dieser Datenbank.

\subsection{Datenbank wechseln}
Im Dropdown oben in \menu{Setup} die gewünschte Datenbank auswählen.
Die Auswahl gilt für die gesamte Browser-Sitzung.

% ==========================================================================
\section{Schüler importieren}

\adminonly

\subsection{CSV-Format}
Eine Zeile pro Schüler, Semikolon-getrennt, \textbf{ohne Kopfzeile}:

\begin{center}
\texttt{Nachname;Vorname;Klasse;TT.MM.JJJJ}
\end{center}

\textbf{Beispiel:}
\begin{verbatim}
Müller;Anna;7a;15.03.2012
Schmidt;Tom;7a;22.07.2011
Weber;Lisa;8b;01.11.2010
\end{verbatim}

\subsection{Import durchführen}
\begin{enumerate}[noitemsep]
  \item \menu{Setup} → \menu{Schülerimport}
  \item CSV-Datei hochladen oder per Drag\,\&\,Drop einwerfen
  \item Ergebnis prüfen: hinzugefügte, aktualisierte und fehlerhafte Einträge
\end{enumerate}

\info{Der Import ist wiederholbar — bestehende Schüler werden aktualisiert,
nicht doppelt angelegt. Schüler, die nicht mehr in der CSV stehen, werden entfernt.}

% ==========================================================================
\section{Zeugnisse exportieren}

\adminonly

\subsection{Vorbereitung}
Vor dem Export sicherstellen (Übersicht nutzen):
\begin{itemize}[noitemsep]
  \item Zeugnistag in \menu{Setup} eingetragen
  \item Alle Schüler haben Niveau und Themenurteile
  \item Alle Zeugnistexte sind geschrieben
  \item Kompetenzauswahl vollständig
\end{itemize}

\subsection{Export durchführen}
\begin{enumerate}[noitemsep]
  \item \menu{Admin} → \menu{Export}
  \item Klasse auswählen
  \item Schüler per Checkbox markieren (oder alle auswählen)
  \item \key{Exportieren} klicken
  \item Fortschrittsanzeige abwarten
\end{enumerate}

\subsection{Speicherort der PDFs}
Die PDFs liegen auf dem Server unter:
\begin{center}
\texttt{\textasciitilde/Zeugnisse/[JJJJ]-[HJ|EJ]/[Klasse]/[Nachname]\_[Vorname].pdf}
\end{center}

\warn{Die PDFs werden auf dem Server gespeichert und müssen von dort
abgerufen, gedruckt oder weitergegeben werden.
Sie werden \textbf{nicht} automatisch versendet.}

% ==========================================================================
\section{Benutzerverwaltung}

\adminonly

\menu{Admin} → \menu{Benutzer}.

\begin{itemize}[noitemsep]
  \item \textbf{Neuen Nutzer anlegen:} Benutzername, Passwort und Rolle eingeben
  \item \textbf{Rollen:} Admin (voller Zugriff) oder Lehrer (eingeschränkt)
  \item \textbf{Nutzer löschen:} über den Löschen-Button in der Benutzerliste
\end{itemize}

\info{Das erste Admin-Konto wird beim Erststart direkt auf der Login-Seite
eingerichtet, wenn noch kein Admin existiert.}

\newpage
% ==========================================================================
% TECHNISCHER ANHANG – FÜR ENTWICKLER
% ==========================================================================

{%
  \centering
  \vspace{2em}
  {\color{gray}\rule{\linewidth}{1pt}}\par\vspace{0.5em}
  {\large\bfseries\color{gray} Technischer Anhang}\par
  \vspace{0.3em}
  {\normalsize\color{gray} Für Entwickler und Vorlagen-Anpasser}\par
  \vspace{0.5em}
  {\color{gray}\rule{\linewidth}{1pt}}\par
  \vspace{2em}
}

% ==========================================================================
\section{Lua-Datenformat}

Für jeden Schüler generiert das Backend eine \texttt{.lua}-Datei mit allen
Bewertungsdaten sowie eine \texttt{.tex}-Datei.
LuaLaTeX lädt die Daten per \texttt{require("studentdata")}.

\subsection{Struktur der Lua-Datei}

\begin{verbatim}
student = {
  first_name      = 'Anna',
  last_name       = 'Müller',
  classRoom       = '7a',
  date_of_birth   = '15.03.2012',
  school_year     = '2025/26',
  part_of_year    = 'Halbjahr',
  report_date     = '31.01.2026',
  personal_text   = '...',   -- LaTeX-formatiert (html_to_latex)
  comment         = '...',
  lb              = false,
  gb              = false,
  absenceDaysTotal           = 3,
  absenceDaysUnauthorized    = 0,
  absenceHoursTotal          = 6,
  absenceHoursUnauthorized   = 0,
  subjects = {
    {
      name   = 'Mathematik',
      level  = 2,            -- Zahl (1–3), String oder 'nicht erteilt'
      topics = {
        {
          title       = 'Arithmetik / Algebra',
          grade       = 2,   -- 1–4, 'ne' oder ''
          competences = {
            { description = 'Ich kann rationale Zahlen...' },
          },
        },
      },
    },
  },
}
return student
\end{verbatim}

\subsection{level-Werte}

\begin{tblr}{
  width=\linewidth,
  colspec={Q[l,wd=4cm] X[l]},
  hlines, vlines,
  cell{1}{1-2}={bg=primary,fg=white,font=\bfseries},
}
level-Wert & Darstellung im PDF \\
1, 2, 3 & „Du hast vorwiegend auf Anforderungsebene N gearbeitet." \\
7, 8, 9 & „bis Klasse N ohne Anforderungsebene" \\
\texttt{'nicht erteilt'} & Freitext-Block mit „nicht erteilt" \\
HTML-String & Freitext (html\_to\_latex konvertiert) \\
\texttt{''} (leer) & Kein Niveau-Block \\
\end{tblr}

\subsection{grade-Werte}

\begin{tblr}{
  width=\linewidth,
  colspec={Q[l,wd=2.5cm] X[l]},
  hlines, vlines,
  cell{1}{1-2}={bg=primary,fg=white,font=\bfseries},
}
grade-Wert & Darstellung \\
1 & sehr gut (erste Box angekreuzt) \\
2 & gut \\
3 & teilweise erfüllt \\
4 & nicht erfüllt \\
\texttt{'ne'} & Colspan: „nicht erteilt" \\
\texttt{''} & Colspan: leer \\
\end{tblr}

\subsection{Vorlagen-Dateien}

\begin{tblr}{
  width=\linewidth,
  colspec={Q[l,wd=5cm] X[l]},
  hlines, vlines,
  cell{1}{1-2}={bg=primary,fg=white,font=\bfseries},
}
Datei & Zweck \\
\texttt{Zeugnis.tex} & Hauptdokument \\
\texttt{ZeugnisPackages.tex} & Pakete, Farben, Befehle \\
\texttt{ZeugnisTitlepage.tex} & Titelseite (Name, Klasse, Datum) \\
\texttt{ZeugnisTables.tex} & Lua-Loop: Fächer → Tabellen \\
\texttt{ZeugnisBack.tex} & Rückseite (Fehltage, Unterschriften) \\
\end{tblr}

\subsection{Textformatierung (html\_to\_latex)}

Der im Editor eingegebene Text wird von \texttt{html\_to\_latex.py} in
tabularray-sicheres LaTeX umgewandelt.
Block-Level-Umgebungen werden vermieden, da sie in tabularray-Zellen
nicht zuverlässig funktionieren.
Listen werden als Hängend-Einzug-Absätze mit \texttt{\textbackslash textbullet} umgesetzt.

\vfill
\begin{center}
\small\color{gray}
Kompetenzen-Tool · Evangelische Gemeinschaftsschule Erfurt ·
\texttt{http://zeugnistool.tgsef.intern}\\
Quellcode und Vorlagen: internes Repository
\end{center}

\end{document}
